\documentclass[../main.tex]{subfiles}
\begin{document}

\chapter{Introduction to Operating Systems}
\lessoninfo{
  video={P1L2},
  textbook={OSTEP \cite{ostep} ch.~2, 6; Silberschatz et al.\ \cite{silberschatz2018} ch.~1--2},
  keywords={operating system, abstraction, arbitration, mechanism, policy, kernel mode, user mode, trap, system call, signal, monolithic operating system, modular operating system, microkernel, Linux},
}

Every application depends on an operating system to use the CPUs, the
memory and the devices of the machine on which it runs.  This lesson defines
what an operating system is and which roles it plays, introduces the
elements from which it is built---abstractions, mechanisms and
policies---and the principles that guide their design.  It then examines
the protection boundary between applications and the operating system,
which applications cross by means of system calls, and compares the ways in
which operating systems are organized, with Linux and Mac OS X as examples.
The abstractions introduced here, such as processes, threads, memory pages
and files, are the subject of the lessons that follow.

\section{What is an operating system?}
\label{sec:l01-what}

A computer system consists of hardware---one or more CPUs, each possibly
with several cores, main memory, storage devices such as hard disks, SSDs
and USB flash drives, network interfaces for Ethernet or WiFi, and graphics
processors (GPUs)---and of the applications that use
it.\source{P1L2, what is an operating system, operating system definition;
OSTEP ch.~2; Silberschatz et al.\ ch.~1.}  Several applications usually run
at the same time, and all of them need these hardware components.  The
operating system is the layer of software that sits between the two
(\cref{fig:l01-os-layers}).

\begin{definition}[Operating system]
  An \term{operating system}[オペレーティングシステム] (OS) is a layer of
  systems software that
  \begin{itemize}
    \item has direct, privileged access to the underlying hardware;
    \item hides the complexity of the hardware from applications;
    \item manages the hardware on behalf of one or more applications
      according to predefined policies; and
    \item ensures that applications are isolated and protected from one
      another.
  \end{itemize}
\end{definition}
\index{OS|see{operating system}}

\begin{figure}[htbp]
  \centering
  % The operating system as the layer between applications and hardware.
% Usage: % The operating system as the layer between applications and hardware.
% Usage: % The operating system as the layer between applications and hardware.
% Usage: \input{figures/l01-os-layers}   (width ~110 mm)
\begin{tikzpicture}[x=1mm, y=1mm, font=\footnotesize\sffamily,
  app/.style={draw, rounded corners=2pt, fill=white, minimum height=7mm,
              minimum width=30mm, inner sep=1pt},
  comp/.style={draw, fill=white, minimum height=7mm, minimum width=23mm,
               inner sep=1pt, align=center, font=\scriptsize\sffamily},
  hw/.style={draw, fill=ln-box, minimum height=8mm, minimum width=16mm,
             inner sep=1pt, align=center, font=\scriptsize\sffamily},
  lab/.style={font=\scriptsize\sffamily, text=ln-muted},
  >={Stealth[length=1.6mm]}]
  % ---- applications
  \foreach \x in {17,55,93} {
    \node[app] at (\x,44) {\iflangja{アプリケーション}{application}};
    \draw[<->] (\x,40.5) -- (\x,32);
  }
  \node[lab, anchor=west] at (94.5,36.2) {\iflangja{システムコール}{system calls}};
  % ---- operating system
  \draw[thick, draw=ln-accent, fill=ln-box, rounded corners=2pt]
    (0,14) rectangle (110,32);
  \node[font=\footnotesize\sffamily\bfseries, text=ln-accent, anchor=north west]
    at (1,31.5) {\iflangja{オペレーティングシステム}{operating system}};
  \node[comp] at (14,20.5)   {\iflangja{スケジューラ}{scheduler}};
  \node[comp] at (41.3,20.5) {\iflangja{メモリマネージャ}{memory manager}};
  \node[comp] at (68.6,20.5) {\iflangja{ファイルシステム}{file systems}};
  \node[comp] at (96,20.5)   {\iflangja{デバイスドライバ}{device drivers}};
  % ---- hardware
  \foreach \x/\t in {8/{CPU},
                     26.8/{\iflangja{メモリ}{memory}},
                     45.6/{GPU},
                     64.4/{\iflangja{ネットワーク\\カード}{network\\card}},
                     83.2/{\iflangja{ディスク，\\SSD}{disk,\\SSD}},
                     102/{\iflangja{USB\\デバイス}{USB\\devices}}} {
    \node[hw] at (\x,4) {\t};
    \draw[<->] (\x,8) -- (\x,14);
  }
\end{tikzpicture}
   (width ~110 mm)
\begin{tikzpicture}[x=1mm, y=1mm, font=\footnotesize\sffamily,
  app/.style={draw, rounded corners=2pt, fill=white, minimum height=7mm,
              minimum width=30mm, inner sep=1pt},
  comp/.style={draw, fill=white, minimum height=7mm, minimum width=23mm,
               inner sep=1pt, align=center, font=\scriptsize\sffamily},
  hw/.style={draw, fill=ln-box, minimum height=8mm, minimum width=16mm,
             inner sep=1pt, align=center, font=\scriptsize\sffamily},
  lab/.style={font=\scriptsize\sffamily, text=ln-muted},
  >={Stealth[length=1.6mm]}]
  % ---- applications
  \foreach \x in {17,55,93} {
    \node[app] at (\x,44) {\iflangja{アプリケーション}{application}};
    \draw[<->] (\x,40.5) -- (\x,32);
  }
  \node[lab, anchor=west] at (94.5,36.2) {\iflangja{システムコール}{system calls}};
  % ---- operating system
  \draw[thick, draw=ln-accent, fill=ln-box, rounded corners=2pt]
    (0,14) rectangle (110,32);
  \node[font=\footnotesize\sffamily\bfseries, text=ln-accent, anchor=north west]
    at (1,31.5) {\iflangja{オペレーティングシステム}{operating system}};
  \node[comp] at (14,20.5)   {\iflangja{スケジューラ}{scheduler}};
  \node[comp] at (41.3,20.5) {\iflangja{メモリマネージャ}{memory manager}};
  \node[comp] at (68.6,20.5) {\iflangja{ファイルシステム}{file systems}};
  \node[comp] at (96,20.5)   {\iflangja{デバイスドライバ}{device drivers}};
  % ---- hardware
  \foreach \x/\t in {8/{CPU},
                     26.8/{\iflangja{メモリ}{memory}},
                     45.6/{GPU},
                     64.4/{\iflangja{ネットワーク\\カード}{network\\card}},
                     83.2/{\iflangja{ディスク，\\SSD}{disk,\\SSD}},
                     102/{\iflangja{USB\\デバイス}{USB\\devices}}} {
    \node[hw] at (\x,4) {\t};
    \draw[<->] (\x,8) -- (\x,14);
  }
\end{tikzpicture}
   (width ~110 mm)
\begin{tikzpicture}[x=1mm, y=1mm, font=\footnotesize\sffamily,
  app/.style={draw, rounded corners=2pt, fill=white, minimum height=7mm,
              minimum width=30mm, inner sep=1pt},
  comp/.style={draw, fill=white, minimum height=7mm, minimum width=23mm,
               inner sep=1pt, align=center, font=\scriptsize\sffamily},
  hw/.style={draw, fill=ln-box, minimum height=8mm, minimum width=16mm,
             inner sep=1pt, align=center, font=\scriptsize\sffamily},
  lab/.style={font=\scriptsize\sffamily, text=ln-muted},
  >={Stealth[length=1.6mm]}]
  % ---- applications
  \foreach \x in {17,55,93} {
    \node[app] at (\x,44) {\iflangja{アプリケーション}{application}};
    \draw[<->] (\x,40.5) -- (\x,32);
  }
  \node[lab, anchor=west] at (94.5,36.2) {\iflangja{システムコール}{system calls}};
  % ---- operating system
  \draw[thick, draw=ln-accent, fill=ln-box, rounded corners=2pt]
    (0,14) rectangle (110,32);
  \node[font=\footnotesize\sffamily\bfseries, text=ln-accent, anchor=north west]
    at (1,31.5) {\iflangja{オペレーティングシステム}{operating system}};
  \node[comp] at (14,20.5)   {\iflangja{スケジューラ}{scheduler}};
  \node[comp] at (41.3,20.5) {\iflangja{メモリマネージャ}{memory manager}};
  \node[comp] at (68.6,20.5) {\iflangja{ファイルシステム}{file systems}};
  \node[comp] at (96,20.5)   {\iflangja{デバイスドライバ}{device drivers}};
  % ---- hardware
  \foreach \x/\t in {8/{CPU},
                     26.8/{\iflangja{メモリ}{memory}},
                     45.6/{GPU},
                     64.4/{\iflangja{ネットワーク\\カード}{network\\card}},
                     83.2/{\iflangja{ディスク，\\SSD}{disk,\\SSD}},
                     102/{\iflangja{USB\\デバイス}{USB\\devices}}} {
    \node[hw] at (\x,4) {\t};
    \draw[<->] (\x,8) -- (\x,14);
  }
\end{tikzpicture}

  \caption{The operating system between applications and hardware.
    Applications request services through system calls; components such
    as the scheduler, the memory manager, file systems and device drivers
    carry them out on the hardware.}
  \label{fig:l01-os-layers}
\end{figure}

In short, an operating system \emph{abstracts} the hardware, simplifying
what it looks like to applications, and \emph{arbitrates} its use,
overseeing and controlling how applications share it.  The definition
names three roles.
\begin{description}
  \item[Hiding hardware complexity] The operating system presents the
    hardware to applications through an
    \term{abstraction}[抽象化]: a simpler, higher-level interface that
    hides the details of the devices behind it.  An application that
    stores data does not address the sectors and blocks of a particular
    disk; it writes to a \emph{file}, and the operating system translates
    file operations into operations on whatever storage device is
    present.  In the same way a \emph{socket} represents network
    communication independently of the network interface.  Because the
    abstraction stays the same when the device changes, applications run
    unchanged on different hardware.
  \item[Resource management] The operating system manages the hardware
    resources by \term{arbitration}[調停]: it decides which application
    uses which resource and how much of it, allocating memory to
    applications, scheduling them on the CPUs and controlling their access
    to devices.  The decisions follow policies, for example that access
    to resources be fair or that the resources used by one application be
    limited.
  \item[Isolation and protection] Applications that run at the same time
    must not interfere with one another.  The operating system provides
    \term{isolation}[隔離] by giving each application its own share of the
    resources, for example its own portion of physical memory, so that its
    execution is not affected by the others.  It provides
    \term{protection}[保護] by preventing an application from accessing
    resources that have not been assigned to it, such as the memory of
    other applications or of the operating system itself.
\end{description}

The operating system thus consists of components such as the CPU scheduler,
the memory manager, device drivers and file systems.  Application programs
that merely use these services, such as a spreadsheet or a media player,
are not part of it, even when they are distributed together with it.

\begin{example}
  A function of the operating system is an abstraction when it simplifies
  the interface to a device, and arbitration when it divides a resource
  among applications.  Letting a program send a document to a printer
  without knowing the command language of the printer's model is
  abstraction; deciding in which order the documents of three programs
  that print at the same time reach the printer is arbitration.  Likewise,
  offering the same socket interface over Ethernet and over WiFi is
  abstraction, while limiting the network bandwidth that one application
  may consume is arbitration.
\end{example}

\subsection{Examples of operating systems}

Operating systems differ with the environment they are designed
for.\source{P1L2, operating system examples.}  On desktop and server
machines the main examples are Microsoft Windows and the UNIX-based
systems, among them Mac OS X, which derives from BSD UNIX, and Linux.  On
embedded devices such as smartphones they include Android, which is based
on Linux, iOS and Symbian.  Each makes design choices suited to its
environment; the lessons that follow use Linux as their main example.

\section{Abstractions, mechanisms and policies}
\label{sec:l01-elements}

An operating system is built from three kinds of
elements.\source{P1L2, OS elements, OS elements example; Silberschatz et
al.\ ch.~2.}  Its abstractions represent executing applications---the
process and the thread---and hardware resources---the file, the socket and
the memory page.  On these abstractions it provides operations, and it
decides how the operations are applied.

\begin{definition}[Mechanism and policy]
  A \term{mechanism}[機構] is an operation that the operating system
  provides on its abstractions, such as creating or scheduling a process,
  opening or writing a file, or allocating memory.  A
  \term{policy}[方針] is a rule that determines how the mechanisms are
  used to manage the hardware, such as least recently used (LRU) for
  memory or earliest deadline first (EDF) for scheduling.
\end{definition}

Memory management shows how the three elements work together
(\cref{tab:l01-elements}).  The abstraction is the memory page, a
fixed-size block of memory.  The mechanisms allocate a page and map it into
the address space of a process, so that the process can use it.  Physical
memory (DRAM) is usually smaller than the total demand of all processes,
so the operating system also moves some pages to disk and brings them back
when they are accessed again.  The policy decides which pages to move: with
LRU, the pages that have not been accessed for the longest time, on the
assumption that they are the least likely to be needed soon.

\begin{table}[htbp]
  \centering
  \caption{The elements of an operating system, in general and for memory
    management.}
  \label{tab:l01-elements}
  \small
  \begin{tabular}{@{}l>{\raggedright}p{47mm}>{\raggedright\arraybackslash}p{45mm}@{}}
    \toprule
    Element & Examples & Memory management \\
    \midrule
    Abstraction & process, thread, file, socket, memory~page & memory page \\
    Mechanism & create, schedule, open, write, allocate
      & allocate a page, map it into a process \\
    Policy & LRU, EDF & LRU: move the least recently used pages to disk \\
    \bottomrule
  \end{tabular}
\end{table}

\section{Design principles}
\label{sec:l01-principles}

Two principles guide the design of these elements.\source{P1L2, OS design
principles; Silberschatz et al.\ ch.~2.}  The first is the
\term{separation of mechanism and policy}[機構と方針の分離]: the operating
system implements flexible mechanisms that can support many policies, so
that a policy can be chosen or replaced without changing the mechanisms.
For choosing the pages to move to disk, the same mechanisms---recording
page accesses, moving a page to disk and updating its mapping---support
LRU, least frequently used (LFU) and even random selection.  Which of them
works best depends on the setting in which the system is used.

\begin{example}
  Physical memory holds four pages, and a fifth page is needed, so one of
  the four must move to disk.  At time \SI{160}{\milli\second} the
  recorded accesses are:
  \begin{center}
    \small
    \begin{tabular}{@{}lcccc@{}}
      \toprule
      Page & A & B & C & D \\
      \midrule
      Last access (\si{\milli\second}) & 120 & 40 & 95 & 150 \\
      Number of accesses & 9 & 25 & 3 & 12 \\
      \bottomrule
    \end{tabular}
  \end{center}
  LRU moves page B, whose last access lies furthest in the past; LFU moves
  page C, which has been accessed least often; random selection moves each
  page with probability $1/4$.  In all three cases the mechanism executed
  is the same: the chosen page is written to disk and its mapping is
  invalidated.
\end{example}

The second principle is to \term{optimize for the common case}[頻繁に起こる場合の最適化].
The designers must understand where the operating system will be used,
which applications its users will run on the machine, and what the
requirements of that workload are.  They then choose the policies that
make sense for this common case and that the underlying abstractions and
mechanisms can support.  A server that mostly runs a database, for
example, performs many random reads and writes, whereas a media server
reads large files sequentially, and the policies that suit one workload
are a poor fit for the other.

\begin{keypoint}
  Abstractions simplify, mechanisms operate on the abstractions, and
  policies decide how the mechanisms are used.  Keeping mechanisms general
  lets the same operating system adopt the policies that fit its common
  case.
\end{keypoint}

\section{The user/kernel protection boundary}
\label{sec:l01-boundary}

Because the operating system needs direct access to the hardware, it must
execute with privileges that applications do not
have.\source{P1L2, user/kernel protection boundary; OSTEP ch.~6;
Silberschatz et al.\ ch.~1.}
Computer platforms therefore distinguish at least two modes of execution.

\begin{definition}[Kernel mode and user mode]
  In \term{kernel mode}[カーネルモード], the privileged mode, software may
  execute every instruction of the processor and access the hardware
  directly.  In \term{user mode}[ユーザモード], the unprivileged mode,
  instructions that access the hardware directly are forbidden.
\end{definition}

The part of the operating system that executes in kernel mode is its
\term{kernel}[カーネル]; applications execute in user mode.  The
distinction between the two modes is supported by the hardware.  The CPU
contains a \term{mode bit}[モードビット] that records the current mode, and
an instruction that directly manipulates hardware---a
\term{privileged instruction}[特権命令]---is executed only if the mode bit
indicates kernel mode.\margin{x86 processors provide four privilege
levels, the protection rings of Lesson~12; operating systems use at least
two of them.}  An attempt to execute a privileged instruction in user mode
does not take effect but causes a trap.

\begin{definition}[Trap]
  A \term{trap}[トラップ] is a transfer of control from a running
  application to the operating system, performed by the hardware: the
  application is interrupted, the CPU switches to kernel mode, and
  execution continues at a location in the kernel that the operating
  system has registered in advance.
\end{definition}

After a trap caused by a privileged instruction or by an access to memory
that requires privileges, the operating system determines what caused the
trap and whether the access should be granted; if the application attempted
something illegal, the operating system may terminate it.  Applications
and the operating system interact in three ways (\cref{fig:l01-user-kernel}):
\begin{description}
  \item[Traps] transfer control to the operating system when an
    application performs an operation that requires privileges.
  \item[System calls] are requested explicitly.  A
    \term{system call}[システムコール] is an operation that the operating
    system exports so that applications can ask it to perform a service or
    a privileged access on their behalf, for example \texttt{open} for a
    file, \texttt{send} for a socket or \texttt{mmap} for memory.  The set
    of these operations is the \emph{system call interface}.
  \item[Signals] travel in the opposite direction.  A
    \term{signal}[シグナル] is a notification that the operating system
    passes to an application, for example that a timer the application
    has set has expired.
\end{description}

\begin{figure}[htbp]
  \centering
  % The user/kernel protection boundary and the ways to cross it.
% Usage: % The user/kernel protection boundary and the ways to cross it.
% Usage: % The user/kernel protection boundary and the ways to cross it.
% Usage: \input{figures/l01-user-kernel}   (width ~112 mm)
\begin{tikzpicture}[x=1mm, y=1mm, font=\footnotesize\sffamily,
  app/.style={draw, rounded corners=2pt, fill=white, minimum height=8mm,
              minimum width=30mm, inner sep=1pt},
  big/.style={draw, rounded corners=2pt, minimum height=8mm,
              minimum width=112mm, inner sep=1pt},
  lab/.style={font=\scriptsize\sffamily, align=center},
  lvl/.style={font=\scriptsize\sffamily, text=ln-muted, align=right},
  >={Stealth[length=1.8mm]}]
  \node[app] (a) at (28,42) {\iflangja{アプリケーション}{application}};
  \node[app] (b) at (78,42) {\iflangja{アプリケーション}{application}};
  \draw[dashed, ln-rule, thick] (0,26) -- (112,26);
  \node[lvl, anchor=south east] at (112,26.8) {\iflangja{ユーザレベル（非特権）}{user level (unprivileged)}};
  \node[lvl, anchor=north east] at (112,25.2) {\iflangja{カーネルレベル（特権）}{kernel level (privileged)}};
  \node[big, fill=ln-box, thick, draw=ln-accent] (k) at (56,15)
    {\iflangja{OS カーネル}{OS kernel}};
  \node[big, fill=white] (h) at (56,0) {\iflangja{ハードウェア}{hardware}};
  % system call and signal
  \draw[->, thick, ln-accent] (22,38) -- (22,19);
  \node[lab, anchor=east] at (21,32) {\iflangja{システムコール}{system call}};
  \draw[->, thick, ln-accent] (34,19) -- (34,38);
  \node[lab, anchor=west] at (35,32) {\iflangja{シグナル}{signal}};
  % trap
  \draw[->, thick, ln-caution] (78,38) -- (78,19);
  \node[lab, anchor=east] at (77,32)
    {\iflangja{トラップ\\（特権操作）}{trap (privileged\\operation)}};
  % privileged access
  \draw[<->, thick] (56,11) -- (56,4);
  \node[lab, anchor=west] at (57,7.5)
    {\iflangja{ハードウェアへの直接アクセス}{direct hardware access}};
\end{tikzpicture}
   (width ~112 mm)
\begin{tikzpicture}[x=1mm, y=1mm, font=\footnotesize\sffamily,
  app/.style={draw, rounded corners=2pt, fill=white, minimum height=8mm,
              minimum width=30mm, inner sep=1pt},
  big/.style={draw, rounded corners=2pt, minimum height=8mm,
              minimum width=112mm, inner sep=1pt},
  lab/.style={font=\scriptsize\sffamily, align=center},
  lvl/.style={font=\scriptsize\sffamily, text=ln-muted, align=right},
  >={Stealth[length=1.8mm]}]
  \node[app] (a) at (28,42) {\iflangja{アプリケーション}{application}};
  \node[app] (b) at (78,42) {\iflangja{アプリケーション}{application}};
  \draw[dashed, ln-rule, thick] (0,26) -- (112,26);
  \node[lvl, anchor=south east] at (112,26.8) {\iflangja{ユーザレベル（非特権）}{user level (unprivileged)}};
  \node[lvl, anchor=north east] at (112,25.2) {\iflangja{カーネルレベル（特権）}{kernel level (privileged)}};
  \node[big, fill=ln-box, thick, draw=ln-accent] (k) at (56,15)
    {\iflangja{OS カーネル}{OS kernel}};
  \node[big, fill=white] (h) at (56,0) {\iflangja{ハードウェア}{hardware}};
  % system call and signal
  \draw[->, thick, ln-accent] (22,38) -- (22,19);
  \node[lab, anchor=east] at (21,32) {\iflangja{システムコール}{system call}};
  \draw[->, thick, ln-accent] (34,19) -- (34,38);
  \node[lab, anchor=west] at (35,32) {\iflangja{シグナル}{signal}};
  % trap
  \draw[->, thick, ln-caution] (78,38) -- (78,19);
  \node[lab, anchor=east] at (77,32)
    {\iflangja{トラップ\\（特権操作）}{trap (privileged\\operation)}};
  % privileged access
  \draw[<->, thick] (56,11) -- (56,4);
  \node[lab, anchor=west] at (57,7.5)
    {\iflangja{ハードウェアへの直接アクセス}{direct hardware access}};
\end{tikzpicture}
   (width ~112 mm)
\begin{tikzpicture}[x=1mm, y=1mm, font=\footnotesize\sffamily,
  app/.style={draw, rounded corners=2pt, fill=white, minimum height=8mm,
              minimum width=30mm, inner sep=1pt},
  big/.style={draw, rounded corners=2pt, minimum height=8mm,
              minimum width=112mm, inner sep=1pt},
  lab/.style={font=\scriptsize\sffamily, align=center},
  lvl/.style={font=\scriptsize\sffamily, text=ln-muted, align=right},
  >={Stealth[length=1.8mm]}]
  \node[app] (a) at (28,42) {\iflangja{アプリケーション}{application}};
  \node[app] (b) at (78,42) {\iflangja{アプリケーション}{application}};
  \draw[dashed, ln-rule, thick] (0,26) -- (112,26);
  \node[lvl, anchor=south east] at (112,26.8) {\iflangja{ユーザレベル（非特権）}{user level (unprivileged)}};
  \node[lvl, anchor=north east] at (112,25.2) {\iflangja{カーネルレベル（特権）}{kernel level (privileged)}};
  \node[big, fill=ln-box, thick, draw=ln-accent] (k) at (56,15)
    {\iflangja{OS カーネル}{OS kernel}};
  \node[big, fill=white] (h) at (56,0) {\iflangja{ハードウェア}{hardware}};
  % system call and signal
  \draw[->, thick, ln-accent] (22,38) -- (22,19);
  \node[lab, anchor=east] at (21,32) {\iflangja{システムコール}{system call}};
  \draw[->, thick, ln-accent] (34,19) -- (34,38);
  \node[lab, anchor=west] at (35,32) {\iflangja{シグナル}{signal}};
  % trap
  \draw[->, thick, ln-caution] (78,38) -- (78,19);
  \node[lab, anchor=east] at (77,32)
    {\iflangja{トラップ\\（特権操作）}{trap (privileged\\operation)}};
  % privileged access
  \draw[<->, thick] (56,11) -- (56,4);
  \node[lab, anchor=west] at (57,7.5)
    {\iflangja{ハードウェアへの直接アクセス}{direct hardware access}};
\end{tikzpicture}

  \caption{The user/kernel protection boundary.  Only the kernel accesses
    the hardware; applications enter the kernel through system calls or
    traps, and the kernel notifies applications with signals.}
  \label{fig:l01-user-kernel}
\end{figure}

\section{System calls}
\label{sec:l01-syscalls}

\Cref{fig:l01-syscall-flow} shows the sequence of events when a user
process executes a system call.\source{P1L2, system call flowchart;
OSTEP ch.~6; Silberschatz et al.\ ch.~1--2.}  The process runs in user mode
until it needs a service of the operating system and calls the system call.
The call traps into the kernel: the mode bit is switched to kernel mode,
the execution context changes from that of the user process to that of the
kernel, the arguments are passed, and execution jumps to the kernel code
that implements the call.  The kernel executes the call in kernel mode.
When it has completed, the mode bit is switched back to user mode, the
execution context of the process is restored, the results are passed back,
and the process continues with the instruction after the call.

\begin{figure}[htbp]
  \centering
  % Sequence of a system call: transition from user mode to kernel mode and back.
% Usage: % Sequence of a system call: transition from user mode to kernel mode and back.
% Usage: % Sequence of a system call: transition from user mode to kernel mode and back.
% Usage: \input{figures/l01-syscall-flow}   (width ~118 mm)
\begin{tikzpicture}[x=1mm, y=1mm, font=\footnotesize\sffamily,
  st/.style={draw, rounded corners=2pt, align=center, minimum height=9mm,
             inner sep=1pt, font=\scriptsize\sffamily},
  lab/.style={font=\scriptsize\sffamily, align=center},
  lvl/.style={font=\scriptsize\sffamily\bfseries, text=ln-muted},
  >={Stealth[length=1.8mm]}]
  % kernel-mode region below the dashed boundary
  \fill[ln-box] (0,-6) rectangle (118,15);
  \draw[dashed, ln-rule, thick] (0,15) -- (118,15);
  \node[lvl, anchor=north west] at (0,32)
    {\iflangja{ユーザモード（モードビット = 1）}{user mode (mode bit = 1)}};
  \node[lvl, anchor=south west] at (0,-5.5)
    {\iflangja{カーネルモード（モードビット = 0）}{kernel mode (mode bit = 0)}};
  \node[st, fill=white, minimum width=24mm, text width=22mm] (a) at (13,23)
    {\iflangja{ユーザプロセス\\実行中}{user process\\executing}};
  \node[st, fill=white, minimum width=22mm, text width=20mm] (b) at (43,23)
    {\iflangja{システムコール\\を呼び出す}{calls\\system call}};
  \node[st, fill=white, thick, draw=ln-accent, minimum width=24mm, text width=22mm]
    (c) at (68,4) {\iflangja{システムコール\\を実行}{execute\\system call}};
  \node[st, fill=white, minimum width=24mm, text width=22mm] (d) at (93,23)
    {\iflangja{システムコール\\から復帰}{return from\\system call}};
  \draw[->] (a) -- (b);
  \draw[->, thick, ln-accent] (b.south) |- (c.west);
  \node[lab, anchor=east] at (42,8.5)
    {\iflangja{トラップ\\モードビット = 0}{trap\\mode bit = 0}};
  \draw[->, thick, ln-accent] (c.east) -| (d.south);
  \node[lab, anchor=west] at (94,8.5)
    {\iflangja{復帰\\モードビット = 1}{return\\mode bit = 1}};
  \draw[->] (d.east) -- (118,23);
\end{tikzpicture}
   (width ~118 mm)
\begin{tikzpicture}[x=1mm, y=1mm, font=\footnotesize\sffamily,
  st/.style={draw, rounded corners=2pt, align=center, minimum height=9mm,
             inner sep=1pt, font=\scriptsize\sffamily},
  lab/.style={font=\scriptsize\sffamily, align=center},
  lvl/.style={font=\scriptsize\sffamily\bfseries, text=ln-muted},
  >={Stealth[length=1.8mm]}]
  % kernel-mode region below the dashed boundary
  \fill[ln-box] (0,-6) rectangle (118,15);
  \draw[dashed, ln-rule, thick] (0,15) -- (118,15);
  \node[lvl, anchor=north west] at (0,32)
    {\iflangja{ユーザモード（モードビット = 1）}{user mode (mode bit = 1)}};
  \node[lvl, anchor=south west] at (0,-5.5)
    {\iflangja{カーネルモード（モードビット = 0）}{kernel mode (mode bit = 0)}};
  \node[st, fill=white, minimum width=24mm, text width=22mm] (a) at (13,23)
    {\iflangja{ユーザプロセス\\実行中}{user process\\executing}};
  \node[st, fill=white, minimum width=22mm, text width=20mm] (b) at (43,23)
    {\iflangja{システムコール\\を呼び出す}{calls\\system call}};
  \node[st, fill=white, thick, draw=ln-accent, minimum width=24mm, text width=22mm]
    (c) at (68,4) {\iflangja{システムコール\\を実行}{execute\\system call}};
  \node[st, fill=white, minimum width=24mm, text width=22mm] (d) at (93,23)
    {\iflangja{システムコール\\から復帰}{return from\\system call}};
  \draw[->] (a) -- (b);
  \draw[->, thick, ln-accent] (b.south) |- (c.west);
  \node[lab, anchor=east] at (42,8.5)
    {\iflangja{トラップ\\モードビット = 0}{trap\\mode bit = 0}};
  \draw[->, thick, ln-accent] (c.east) -| (d.south);
  \node[lab, anchor=west] at (94,8.5)
    {\iflangja{復帰\\モードビット = 1}{return\\mode bit = 1}};
  \draw[->] (d.east) -- (118,23);
\end{tikzpicture}
   (width ~118 mm)
\begin{tikzpicture}[x=1mm, y=1mm, font=\footnotesize\sffamily,
  st/.style={draw, rounded corners=2pt, align=center, minimum height=9mm,
             inner sep=1pt, font=\scriptsize\sffamily},
  lab/.style={font=\scriptsize\sffamily, align=center},
  lvl/.style={font=\scriptsize\sffamily\bfseries, text=ln-muted},
  >={Stealth[length=1.8mm]}]
  % kernel-mode region below the dashed boundary
  \fill[ln-box] (0,-6) rectangle (118,15);
  \draw[dashed, ln-rule, thick] (0,15) -- (118,15);
  \node[lvl, anchor=north west] at (0,32)
    {\iflangja{ユーザモード（モードビット = 1）}{user mode (mode bit = 1)}};
  \node[lvl, anchor=south west] at (0,-5.5)
    {\iflangja{カーネルモード（モードビット = 0）}{kernel mode (mode bit = 0)}};
  \node[st, fill=white, minimum width=24mm, text width=22mm] (a) at (13,23)
    {\iflangja{ユーザプロセス\\実行中}{user process\\executing}};
  \node[st, fill=white, minimum width=22mm, text width=20mm] (b) at (43,23)
    {\iflangja{システムコール\\を呼び出す}{calls\\system call}};
  \node[st, fill=white, thick, draw=ln-accent, minimum width=24mm, text width=22mm]
    (c) at (68,4) {\iflangja{システムコール\\を実行}{execute\\system call}};
  \node[st, fill=white, minimum width=24mm, text width=22mm] (d) at (93,23)
    {\iflangja{システムコール\\から復帰}{return from\\system call}};
  \draw[->] (a) -- (b);
  \draw[->, thick, ln-accent] (b.south) |- (c.west);
  \node[lab, anchor=east] at (42,8.5)
    {\iflangja{トラップ\\モードビット = 0}{trap\\mode bit = 0}};
  \draw[->, thick, ln-accent] (c.east) -| (d.south);
  \node[lab, anchor=west] at (94,8.5)
    {\iflangja{復帰\\モードビット = 1}{return\\mode bit = 1}};
  \draw[->] (d.east) -- (118,23);
\end{tikzpicture}

  \caption{Execution of a system call: a trap switches to kernel mode, the
    kernel executes the call, and the return switches back to user mode.
    The mode bit follows the convention of Silberschatz et al.\
    \cite{silberschatz2018}: 0 in kernel mode, 1 in user mode.}
  \label{fig:l01-syscall-flow}
\end{figure}

To make a system call, an application writes the arguments and other
relevant data to a well-defined location and invokes the call by its
\emph{system call number}.  From the number the kernel determines which
operation is requested, how many arguments it takes and where to find
them.  Arguments are passed either directly, as values, or indirectly, as
the address of a memory area that holds them.  In the synchronous mode
described here the process waits until the system call has completed;
asynchronous system calls are discussed in later lessons.

\begin{example}
  Both calls in the following C program ask Linux to write six bytes to
  the standard output:
  \begin{lstlisting}[language=C, numbers=none]
#define _GNU_SOURCE
#include <unistd.h>
#include <sys/syscall.h>

int main(void) {
  const char msg[] = "hello\n";
  /* C library wrapper for the system call */
  write(1, msg, sizeof msg - 1);
  /* the same system call, invoked by number */
  syscall(SYS_write, 1, msg, sizeof msg - 1);
  return 0;
}
\end{lstlisting}
  The library function \texttt{write} places the number of the
  \texttt{write} system call and its three arguments where the kernel
  expects them and executes the trap.  The file descriptor~1 (standard
  output) and the length are passed directly; the text is passed indirectly, through its address
  \texttt{msg}.\margin{x86-64 Linux: number in \texttt{rax}, arguments in
  \texttt{rdi}, \texttt{rsi}, \texttt{rdx}; the \texttt{syscall}
  instruction traps.}
\end{example}

\subsection{Crossing the user/kernel boundary}

Transitions between user mode and kernel mode are a necessary part of
application execution: an application that reads a file, sends data over
the network or allocates memory must cross the
boundary.\source{P1L2, crossing the OS boundary.}  The transitions have
three properties.
\begin{itemize}
  \item They are supported by the hardware, through traps on illegal
    instructions or on memory accesses that require special privileges.
  \item They require executing a number of instructions to save and
    switch state; on a \SI{2}{\giga\hertz} machine running Linux a
    transition takes about \SIrange{50}{100}{\nano\second}.
  \item They switch locality.  The kernel executes different code and uses
    different data than the application, so its instructions replace
    application data in the processor caches.  When the application
    resumes, part of its data must be fetched from main memory again,
    which slows it down.
\end{itemize}
User/kernel transitions are therefore not cheap.\caution{A system call
looks like a procedure call in C but costs far more.}

\section{Operating system services}
\label{sec:l01-services}

An operating system gives applications access to the hardware by exporting
services.\source{P1L2, OS services; Silberschatz et al.\ ch.~2.}  At the
most basic level the services correspond to hardware components: a
scheduler controls access to the CPUs; a memory manager allocates physical
memory to the applications that run at the same time and makes sure that
they do not overwrite each other's memory; a block device driver gives
access to block devices such as disks.  Other services provide higher-level
abstractions that do not correspond directly to one hardware component,
such as the files and directories of a file system.  Grouped by purpose,
operating systems offer process management, file management, device
management, memory management, storage management and security.

Applications reach all of these services through system calls.  Different
operating systems name and structure their calls differently, but they
provide comparable functionality, as the Windows and Unix calls in
\cref{tab:l01-syscalls} show.

\begin{table}[htbp]
  \centering
  \caption{Examples of Windows and Unix system calls
    (after Silberschatz et al.\ \cite{silberschatz2018}).}
  \label{tab:l01-syscalls}
  \small
  \begin{tabular}{@{}>{\raggedright}p{24mm}>{\raggedright}p{55mm}%
      >{\raggedright\arraybackslash}p{32mm}@{}}
    \toprule
    Category & Windows & Unix \\
    \midrule
    Process control
      & \texttt{CreateProcess()}, \texttt{ExitProcess()},
        \texttt{WaitForSingleObject()}
      & \texttt{fork()}, \texttt{exit()}, \texttt{wait()} \\
    \addlinespace
    File manipulation
      & \texttt{CreateFile()}, \texttt{ReadFile()}, \texttt{WriteFile()},
        \texttt{CloseHandle()}
      & \texttt{open()}, \texttt{read()}, \texttt{write()},
        \texttt{close()} \\
    \addlinespace
    Device manipulation
      & \texttt{SetConsoleMode()}, \texttt{ReadConsole()},
        \texttt{WriteConsole()}
      & \texttt{ioctl()}, \texttt{read()}, \texttt{write()} \\
    \addlinespace
    Information maintenance
      & \texttt{GetCurrentProcessId()}, \texttt{SetTimer()},
        \texttt{Sleep()}
      & \texttt{getpid()}, \texttt{alarm()}, \texttt{sleep()} \\
    \addlinespace
    Communication
      & \texttt{CreatePipe()}, \texttt{CreateFileMapping()},
        \texttt{MapViewOfFile()}
      & \texttt{pipe()}, \texttt{shm\_open()}, \texttt{mmap()} \\
    \addlinespace
    Protection
      & \texttt{SetFileSecurity()}, \texttt{InitializeSecurityDescriptor()},
        \texttt{SetSecurityDescriptorGroup()}
      & \texttt{chmod()}, \texttt{umask()}, \texttt{chown()} \\
    \bottomrule
  \end{tabular}
\end{table}

\section{Organization of an operating system}
\label{sec:l01-organization}

Operating systems differ in how their services are
organized.\source{P1L2, monolithic OS, modular OS, microkernel;
Silberschatz et al.\ ch.~2.}  \Cref{fig:l01-structures} shows three basic
organizations, and \cref{tab:l01-organization} summarizes their advantages
and disadvantages.

\begin{figure}[htbp]
  \centering
  % Monolithic, modular and microkernel organization of an operating system.
% Usage: % Monolithic, modular and microkernel organization of an operating system.
% Usage: % Monolithic, modular and microkernel organization of an operating system.
% Usage: \input{figures/l01-structures}   (width ~116 mm)
\begin{tikzpicture}[x=1mm, y=1mm, font=\scriptsize\sffamily,
  app/.style={draw, rounded corners=1pt, fill=white, minimum height=5mm,
              minimum width=10mm, inner sep=0.5pt},
  row/.style={draw, fill=white, minimum height=4.6mm, minimum width=32mm,
              inner sep=0.5pt},
  mod/.style={draw, fill=white, minimum height=8mm, minimum width=10mm,
              inner sep=0.5pt, align=center},
  srv/.style={draw, rounded corners=1pt, fill=white, minimum height=8mm,
              minimum width=10.6mm, inner sep=0.5pt, align=center},
  kern/.style={draw=ln-accent, thick, fill=ln-box, rounded corners=2pt},
  hw/.style={draw, fill=white, minimum height=5mm, minimum width=36mm,
             inner sep=0.5pt},
  bnd/.style={dashed, ln-rule, thick},
  ttl/.style={font=\footnotesize\sffamily\bfseries, anchor=south},
  lab/.style={font=\scriptsize\sffamily, text=ln-muted},
  >={Stealth[length=1.4mm]}]
  % ================= monolithic
  \node[ttl] at (18,52) {\iflangja{モノリシック}{monolithic}};
  \foreach \x in {6,18,30} \node[app] at (\x,45) {\iflangja{アプリ}{app}};
  \draw[bnd] (0,39.5) -- (36,39.5);
  \draw[kern] (0,7) rectangle (36,38);
  \foreach \y/\t in {35.25/{\iflangja{プロセス・スレッド}{processes, threads}},
                     30.15/{\iflangja{スケジューリング}{scheduling}},
                     25.05/{\iflangja{メモリ管理}{memory management}},
                     19.95/{\iflangja{デバイスドライバ}{device drivers}},
                     14.85/{\iflangja{FS（逐次アクセス用）}{file system (sequential)}},
                     9.75/{\iflangja{FS（ランダムアクセス用）}{file system (random)}}} {
    \node[row] at (18,\y) {\t};
  }
  \node[hw] at (18,2.5) {\iflangja{ハードウェア}{hardware}};
  % ================= modular
  \begin{scope}[xshift=40mm]
  \node[ttl] at (18,52) {\iflangja{モジュール型}{modular}};
  \foreach \x in {6,18,30} \node[app] at (\x,45) {\iflangja{アプリ}{app}};
  \draw[bnd] (0,39.5) -- (36,39.5);
  \draw[kern] (0,7) rectangle (36,38);
  \node[row, minimum height=9mm, align=center] at (18,31.5)
    {\iflangja{基本サービス\\と API}{basic services\\and APIs}};
  \draw[very thick] (1.5,25.5) -- (34.5,25.5);
  \foreach \x/\t in {7/{FS},
                     18/{\iflangja{ドライバ}{driver}},
                     29/{\iflangja{プロト\\コル}{network\\protocol}}} {
    \node[mod] at (\x,19) {\t};
    \draw[very thick] (\x,23) -- (\x,25.5);
  }
  \node[lab, align=center] at (18,10.5)
    {\iflangja{モジュールを動的に追加}{modules loaded\\dynamically}};
  \node[hw] at (18,2.5) {\iflangja{ハードウェア}{hardware}};
  \end{scope}
  % ================= microkernel
  \begin{scope}[xshift=80mm]
  \node[ttl] at (18,52) {\iflangja{マイクロカーネル}{microkernel}};
  \foreach \x in {6,18,30} \node[app] at (\x,45) {\iflangja{アプリ}{app}};
  \node[srv] (fs) at (6,33)  {\iflangja{ファイル\\システム}{file\\system}};
  \node[srv] (dd) at (18,33) {\iflangja{デバイス\\ドライバ}{device\\driver}};
  \node[srv] (db) at (30,33) {\iflangja{データ\\ベース}{database}};
  \draw[bnd] (0,19.5) -- (36,19.5);
  \draw[kern] (0,7) rectangle (36,18);
  \node[align=center] at (18,12.5)
    {\iflangja{アドレス空間，\\スレッド，IPC}{address spaces,\\threads, IPC}};
  \foreach \n in {fs,dd,db} \draw[<->, ln-accent] (\n.south) -- ++(0,-10);
  \node[hw] at (18,2.5) {\iflangja{ハードウェア}{hardware}};
  \end{scope}
\end{tikzpicture}
   (width ~116 mm)
\begin{tikzpicture}[x=1mm, y=1mm, font=\scriptsize\sffamily,
  app/.style={draw, rounded corners=1pt, fill=white, minimum height=5mm,
              minimum width=10mm, inner sep=0.5pt},
  row/.style={draw, fill=white, minimum height=4.6mm, minimum width=32mm,
              inner sep=0.5pt},
  mod/.style={draw, fill=white, minimum height=8mm, minimum width=10mm,
              inner sep=0.5pt, align=center},
  srv/.style={draw, rounded corners=1pt, fill=white, minimum height=8mm,
              minimum width=10.6mm, inner sep=0.5pt, align=center},
  kern/.style={draw=ln-accent, thick, fill=ln-box, rounded corners=2pt},
  hw/.style={draw, fill=white, minimum height=5mm, minimum width=36mm,
             inner sep=0.5pt},
  bnd/.style={dashed, ln-rule, thick},
  ttl/.style={font=\footnotesize\sffamily\bfseries, anchor=south},
  lab/.style={font=\scriptsize\sffamily, text=ln-muted},
  >={Stealth[length=1.4mm]}]
  % ================= monolithic
  \node[ttl] at (18,52) {\iflangja{モノリシック}{monolithic}};
  \foreach \x in {6,18,30} \node[app] at (\x,45) {\iflangja{アプリ}{app}};
  \draw[bnd] (0,39.5) -- (36,39.5);
  \draw[kern] (0,7) rectangle (36,38);
  \foreach \y/\t in {35.25/{\iflangja{プロセス・スレッド}{processes, threads}},
                     30.15/{\iflangja{スケジューリング}{scheduling}},
                     25.05/{\iflangja{メモリ管理}{memory management}},
                     19.95/{\iflangja{デバイスドライバ}{device drivers}},
                     14.85/{\iflangja{FS（逐次アクセス用）}{file system (sequential)}},
                     9.75/{\iflangja{FS（ランダムアクセス用）}{file system (random)}}} {
    \node[row] at (18,\y) {\t};
  }
  \node[hw] at (18,2.5) {\iflangja{ハードウェア}{hardware}};
  % ================= modular
  \begin{scope}[xshift=40mm]
  \node[ttl] at (18,52) {\iflangja{モジュール型}{modular}};
  \foreach \x in {6,18,30} \node[app] at (\x,45) {\iflangja{アプリ}{app}};
  \draw[bnd] (0,39.5) -- (36,39.5);
  \draw[kern] (0,7) rectangle (36,38);
  \node[row, minimum height=9mm, align=center] at (18,31.5)
    {\iflangja{基本サービス\\と API}{basic services\\and APIs}};
  \draw[very thick] (1.5,25.5) -- (34.5,25.5);
  \foreach \x/\t in {7/{FS},
                     18/{\iflangja{ドライバ}{driver}},
                     29/{\iflangja{プロト\\コル}{network\\protocol}}} {
    \node[mod] at (\x,19) {\t};
    \draw[very thick] (\x,23) -- (\x,25.5);
  }
  \node[lab, align=center] at (18,10.5)
    {\iflangja{モジュールを動的に追加}{modules loaded\\dynamically}};
  \node[hw] at (18,2.5) {\iflangja{ハードウェア}{hardware}};
  \end{scope}
  % ================= microkernel
  \begin{scope}[xshift=80mm]
  \node[ttl] at (18,52) {\iflangja{マイクロカーネル}{microkernel}};
  \foreach \x in {6,18,30} \node[app] at (\x,45) {\iflangja{アプリ}{app}};
  \node[srv] (fs) at (6,33)  {\iflangja{ファイル\\システム}{file\\system}};
  \node[srv] (dd) at (18,33) {\iflangja{デバイス\\ドライバ}{device\\driver}};
  \node[srv] (db) at (30,33) {\iflangja{データ\\ベース}{database}};
  \draw[bnd] (0,19.5) -- (36,19.5);
  \draw[kern] (0,7) rectangle (36,18);
  \node[align=center] at (18,12.5)
    {\iflangja{アドレス空間，\\スレッド，IPC}{address spaces,\\threads, IPC}};
  \foreach \n in {fs,dd,db} \draw[<->, ln-accent] (\n.south) -- ++(0,-10);
  \node[hw] at (18,2.5) {\iflangja{ハードウェア}{hardware}};
  \end{scope}
\end{tikzpicture}
   (width ~116 mm)
\begin{tikzpicture}[x=1mm, y=1mm, font=\scriptsize\sffamily,
  app/.style={draw, rounded corners=1pt, fill=white, minimum height=5mm,
              minimum width=10mm, inner sep=0.5pt},
  row/.style={draw, fill=white, minimum height=4.6mm, minimum width=32mm,
              inner sep=0.5pt},
  mod/.style={draw, fill=white, minimum height=8mm, minimum width=10mm,
              inner sep=0.5pt, align=center},
  srv/.style={draw, rounded corners=1pt, fill=white, minimum height=8mm,
              minimum width=10.6mm, inner sep=0.5pt, align=center},
  kern/.style={draw=ln-accent, thick, fill=ln-box, rounded corners=2pt},
  hw/.style={draw, fill=white, minimum height=5mm, minimum width=36mm,
             inner sep=0.5pt},
  bnd/.style={dashed, ln-rule, thick},
  ttl/.style={font=\footnotesize\sffamily\bfseries, anchor=south},
  lab/.style={font=\scriptsize\sffamily, text=ln-muted},
  >={Stealth[length=1.4mm]}]
  % ================= monolithic
  \node[ttl] at (18,52) {\iflangja{モノリシック}{monolithic}};
  \foreach \x in {6,18,30} \node[app] at (\x,45) {\iflangja{アプリ}{app}};
  \draw[bnd] (0,39.5) -- (36,39.5);
  \draw[kern] (0,7) rectangle (36,38);
  \foreach \y/\t in {35.25/{\iflangja{プロセス・スレッド}{processes, threads}},
                     30.15/{\iflangja{スケジューリング}{scheduling}},
                     25.05/{\iflangja{メモリ管理}{memory management}},
                     19.95/{\iflangja{デバイスドライバ}{device drivers}},
                     14.85/{\iflangja{FS（逐次アクセス用）}{file system (sequential)}},
                     9.75/{\iflangja{FS（ランダムアクセス用）}{file system (random)}}} {
    \node[row] at (18,\y) {\t};
  }
  \node[hw] at (18,2.5) {\iflangja{ハードウェア}{hardware}};
  % ================= modular
  \begin{scope}[xshift=40mm]
  \node[ttl] at (18,52) {\iflangja{モジュール型}{modular}};
  \foreach \x in {6,18,30} \node[app] at (\x,45) {\iflangja{アプリ}{app}};
  \draw[bnd] (0,39.5) -- (36,39.5);
  \draw[kern] (0,7) rectangle (36,38);
  \node[row, minimum height=9mm, align=center] at (18,31.5)
    {\iflangja{基本サービス\\と API}{basic services\\and APIs}};
  \draw[very thick] (1.5,25.5) -- (34.5,25.5);
  \foreach \x/\t in {7/{FS},
                     18/{\iflangja{ドライバ}{driver}},
                     29/{\iflangja{プロト\\コル}{network\\protocol}}} {
    \node[mod] at (\x,19) {\t};
    \draw[very thick] (\x,23) -- (\x,25.5);
  }
  \node[lab, align=center] at (18,10.5)
    {\iflangja{モジュールを動的に追加}{modules loaded\\dynamically}};
  \node[hw] at (18,2.5) {\iflangja{ハードウェア}{hardware}};
  \end{scope}
  % ================= microkernel
  \begin{scope}[xshift=80mm]
  \node[ttl] at (18,52) {\iflangja{マイクロカーネル}{microkernel}};
  \foreach \x in {6,18,30} \node[app] at (\x,45) {\iflangja{アプリ}{app}};
  \node[srv] (fs) at (6,33)  {\iflangja{ファイル\\システム}{file\\system}};
  \node[srv] (dd) at (18,33) {\iflangja{デバイス\\ドライバ}{device\\driver}};
  \node[srv] (db) at (30,33) {\iflangja{データ\\ベース}{database}};
  \draw[bnd] (0,19.5) -- (36,19.5);
  \draw[kern] (0,7) rectangle (36,18);
  \node[align=center] at (18,12.5)
    {\iflangja{アドレス空間，\\スレッド，IPC}{address spaces,\\threads, IPC}};
  \foreach \n in {fs,dd,db} \draw[<->, ln-accent] (\n.south) -- ++(0,-10);
  \node[hw] at (18,2.5) {\iflangja{ハードウェア}{hardware}};
  \end{scope}
\end{tikzpicture}

  \caption{Three organizations of an operating system.  The dashed line is
    the user/kernel boundary.  A monolithic kernel contains every service;
    a modular kernel adds services as modules through a common interface;
    a microkernel keeps only address spaces, threads and IPC, and runs
    all other services at user level.}
  \label{fig:l01-structures}
\end{figure}

\subsection{Monolithic operating systems}

Historically, operating systems were monolithic.  In a
\term{monolithic operating system}[モノリシックOS] every service that any
application may require and that any type of hardware may demand is part
of the operating system: memory management, device drivers, file
management, processes and threads, scheduling.  Such a system may even
include several file systems, one specialized for sequential access and
another optimized for random access, as used by databases.  Because all
of this code is compiled together, the compiler can optimize across the
boundaries between components, for example by inlining.  The price is a
large amount of code and state that is hard to customize, port, maintain,
debug and upgrade, and a large memory footprint that reduces the memory
left for applications and can therefore degrade their performance.

\subsection{Modular operating systems}

A \term{modular operating system}[モジュール型OS], of which Linux is the
standard example, contains a set of basic services and APIs, and
everything else can be added as a \term{kernel module}[カーネルモジュール].
The operating system specifies interfaces that a module must implement to
become part of it, for example an interface for file systems.  Modules can
be installed dynamically, according to the workload: a machine that runs a
database can load a file system optimized for random access.  Since only
the modules that are needed are loaded, the system is smaller, easier to
maintain and upgrade, and leaves more resources to applications.  The
indirection through the module interface can cost some performance,
although usually not much, and modules may come from entirely different
code bases, which makes them a source of bugs.\margin{Silberschatz et al.\
describe Linux as a monolithic kernel with loadable modules.}

\subsection{Microkernels}

A \term{microkernel}[マイクロカーネル] provides only the most basic
primitives at the operating-system level: address spaces, which represent
the memory of executing applications, and threads, which represent their
execution contexts.  Everything else---file systems, device drivers,
databases---runs outside the kernel, at user level, as separate processes.
Serving a request therefore involves many interactions between processes,
and inter-process communication (IPC) is one of the core abstractions and
mechanisms of the microkernel, together with address spaces and threads.
A microkernel is small, so the kernel itself has a small memory footprint
and, for some applications, lower overhead and better performance; its
code is also easier to verify.  This
makes microkernels attractive for embedded devices and control systems,
where correct behavior of the operating system is critical.  On the
other hand, microkernels are often specialized to the underlying hardware,
which limits portability; software development is more complex, because
common operating-system components are harder to find; and the frequent
crossings of the user/kernel boundary are costly.

\begin{table}[htbp]
  \centering
  \caption{Advantages and disadvantages of the three organizations.}
  \label{tab:l01-organization}
  \small
  \begin{tabular}{@{}>{\RaggedRight}p{19mm}>{\RaggedRight}p{46mm}%
      >{\RaggedRight\arraybackslash}p{47mm}@{}}
    \toprule
    Organization & Advantages & Disadvantages \\
    \midrule
    Monolithic & everything included; optimization across components
      & hard to customize, port and manage; large memory footprint \\
    \addlinespace
    Modular & maintainable; smaller footprint; fewer resources needed
      & indirection through interfaces; bugs from separately developed
      modules \\
    \addlinespace
    Microkernel & small kernel footprint; easier to verify
      & limited portability; complex software development; frequent
      user/kernel crossings \\
    \bottomrule
  \end{tabular}
\end{table}

\begin{example}
  Consider a read from a file on disk.  In a monolithic or modular kernel
  the file system and the disk driver run inside the kernel, so the
  application makes one system call and the request crosses the user/kernel
  boundary twice, on entry and on return.  In a microkernel the application
  sends an IPC message to the file-system process, which sends one to the
  driver process; the driver replies to the file system, which replies to
  the application.  Each of the four messages enters the kernel and leaves
  it again, for eight crossings in total.  Assuming \SI{80}{\nano\second}
  per crossing and \num{50000} reads per second, the monolithic design
  spends $\num{50000}\times 2\times\SI{80}{\nano\second} =
  \SI{8}{\milli\second}$ of every second on crossings and the microkernel
  $\num{50000}\times 8\times\SI{80}{\nano\second} =
  \SI{32}{\milli\second}$, not counting the switches between the address
  spaces of the processes involved.
\end{example}

\begin{keypoint}
  A monolithic kernel is complete and can be optimized as a whole but is
  large and hard to manage; a modular kernel adds only the services that
  are needed, at the cost of indirection; a microkernel is small and
  verifiable but pays for frequent user/kernel crossings and IPC.
\end{keypoint}

\section{Linux and Mac OS X}
\label{sec:l01-linux-macos}

A Linux system is organized in layers (\cref{fig:l01-linux},
left).\source{P1L2, Linux and Mac OS architectures; Tanenbaum and Bos
\cite{tanenbaum2015} ch.~10; Silberschatz et al.\ ch.~2.}  The Linux kernel
runs in kernel mode directly on the hardware.  Above it, in user mode, the
standard library provides the system call interface as functions such as
\texttt{open}, \texttt{close}, \texttt{read}, \texttt{write} and
\texttt{fork}; the standard utility programs, such as the shell, editors
and compilers, use the library; and users interact with the utility
programs and with their applications.  The boundaries between these layers
are the system call interface, the library interface and the user
interface.

\begin{figure}[htbp]
  \centering
  % Linux: the layers of a Linux system (left) and the structure of the kernel (right).
% Usage: % Linux: the layers of a Linux system (left) and the structure of the kernel (right).
% Usage: % Linux: the layers of a Linux system (left) and the structure of the kernel (right).
% Usage: \input{figures/l01-linux}   (width ~118 mm)
\begin{tikzpicture}[x=1mm, y=1mm, font=\scriptsize\sffamily,
  lay/.style={draw, rounded corners=1pt, fill=white, minimum width=40mm,
              text width=38mm, align=center, inner sep=0.5pt},
  bar/.style={draw, fill=ln-box, minimum width=72mm, minimum height=5mm,
              inner sep=0.5pt},
  hd/.style={draw, fill=ln-box, minimum width=23.2mm, minimum height=5mm,
             inner sep=0.5pt, font=\scriptsize\sffamily\bfseries},
  cell/.style={draw, fill=white, minimum width=23.2mm, text width=22.2mm,
               minimum height=9mm, align=center, inner sep=0.5pt},
  lab/.style={font=\scriptsize\sffamily, text=ln-muted},
  ttl/.style={font=\footnotesize\sffamily\bfseries, anchor=base}]
  % ---------------- layers
  \node[ttl] at (20,54) {\iflangja{Linux システムの階層}{layers of a Linux system}};
  \node[lay, minimum height=6mm] at (20,48.5) {\iflangja{ユーザ}{users}};
  \node[lay, minimum height=10mm] at (20,39)
    {\iflangja{標準ユーティリティ\\（シェル，エディタ，コンパイラ，…）}{standard utility programs\\(shell, editors, compilers, \dots)}};
  \node[lay, minimum height=10mm] at (20,27.5)
    {\iflangja{標準ライブラリ\\（open, close, read, write, fork, …）}{standard library\\(open, close, read, write, fork, \dots)}};
  \node[lay, minimum height=8mm, fill=ln-box, draw=ln-accent, thick] at (20,15.5)
    {\iflangja{Linux カーネル}{Linux kernel}};
  \node[lay, minimum height=8mm] at (20,5.5)
    {\iflangja{ハードウェア}{hardware}};
  \draw[dashed, ln-rule, thick] (-1,21) -- (41,21);
  % ---------------- kernel structure
  \begin{scope}[xshift=46mm]
  \node[ttl] at (36,54) {\iflangja{Linux カーネルの構成}{structure of the Linux kernel}};
  \node[bar] at (36,46) {\iflangja{システムコール}{system calls}};
  \node[hd] at (11.6,39.5) {I/O};
  \node[hd] at (36,39.5)   {\iflangja{メモリ管理}{memory mgmt.}};
  \node[hd] at (60.4,39.5) {\iflangja{プロセス管理}{process mgmt.}};
  \foreach \y/\a/\b/\c in {
      31.5/{\iflangja{端末\\文字デバイスドライバ}{terminals, character\\device drivers}}/{\iflangja{仮想メモリ}{virtual memory}}/{\iflangja{シグナル処理}{signal handling}},
      21.5/{\iflangja{ソケット，プロトコル，\\ネットワークドライバ}{sockets, protocols,\\network drivers}}/{\iflangja{ページング，\\ページ置換}{paging, page\\replacement}}/{\iflangja{プロセス・スレッドの\\生成と終了}{process and thread\\creation, termination}},
      11.5/{\iflangja{ファイルシステム，\\ブロックデバイスドライバ}{file systems, block\\device drivers}}/{\iflangja{ページキャッシュ}{page cache}}/{\iflangja{CPU スケジューリング}{CPU scheduling}}} {
    \node[cell] at (11.6,\y) {\a};
    \node[cell] at (36,\y)   {\b};
    \node[cell] at (60.4,\y) {\c};
  }
  \node[bar] at (36,3.5) {\iflangja{割り込み，ディスパッチャ}{interrupts, dispatcher}};
  \end{scope}
\end{tikzpicture}
   (width ~118 mm)
\begin{tikzpicture}[x=1mm, y=1mm, font=\scriptsize\sffamily,
  lay/.style={draw, rounded corners=1pt, fill=white, minimum width=40mm,
              text width=38mm, align=center, inner sep=0.5pt},
  bar/.style={draw, fill=ln-box, minimum width=72mm, minimum height=5mm,
              inner sep=0.5pt},
  hd/.style={draw, fill=ln-box, minimum width=23.2mm, minimum height=5mm,
             inner sep=0.5pt, font=\scriptsize\sffamily\bfseries},
  cell/.style={draw, fill=white, minimum width=23.2mm, text width=22.2mm,
               minimum height=9mm, align=center, inner sep=0.5pt},
  lab/.style={font=\scriptsize\sffamily, text=ln-muted},
  ttl/.style={font=\footnotesize\sffamily\bfseries, anchor=base}]
  % ---------------- layers
  \node[ttl] at (20,54) {\iflangja{Linux システムの階層}{layers of a Linux system}};
  \node[lay, minimum height=6mm] at (20,48.5) {\iflangja{ユーザ}{users}};
  \node[lay, minimum height=10mm] at (20,39)
    {\iflangja{標準ユーティリティ\\（シェル，エディタ，コンパイラ，…）}{standard utility programs\\(shell, editors, compilers, \dots)}};
  \node[lay, minimum height=10mm] at (20,27.5)
    {\iflangja{標準ライブラリ\\（open, close, read, write, fork, …）}{standard library\\(open, close, read, write, fork, \dots)}};
  \node[lay, minimum height=8mm, fill=ln-box, draw=ln-accent, thick] at (20,15.5)
    {\iflangja{Linux カーネル}{Linux kernel}};
  \node[lay, minimum height=8mm] at (20,5.5)
    {\iflangja{ハードウェア}{hardware}};
  \draw[dashed, ln-rule, thick] (-1,21) -- (41,21);
  % ---------------- kernel structure
  \begin{scope}[xshift=46mm]
  \node[ttl] at (36,54) {\iflangja{Linux カーネルの構成}{structure of the Linux kernel}};
  \node[bar] at (36,46) {\iflangja{システムコール}{system calls}};
  \node[hd] at (11.6,39.5) {I/O};
  \node[hd] at (36,39.5)   {\iflangja{メモリ管理}{memory mgmt.}};
  \node[hd] at (60.4,39.5) {\iflangja{プロセス管理}{process mgmt.}};
  \foreach \y/\a/\b/\c in {
      31.5/{\iflangja{端末\\文字デバイスドライバ}{terminals, character\\device drivers}}/{\iflangja{仮想メモリ}{virtual memory}}/{\iflangja{シグナル処理}{signal handling}},
      21.5/{\iflangja{ソケット，プロトコル，\\ネットワークドライバ}{sockets, protocols,\\network drivers}}/{\iflangja{ページング，\\ページ置換}{paging, page\\replacement}}/{\iflangja{プロセス・スレッドの\\生成と終了}{process and thread\\creation, termination}},
      11.5/{\iflangja{ファイルシステム，\\ブロックデバイスドライバ}{file systems, block\\device drivers}}/{\iflangja{ページキャッシュ}{page cache}}/{\iflangja{CPU スケジューリング}{CPU scheduling}}} {
    \node[cell] at (11.6,\y) {\a};
    \node[cell] at (36,\y)   {\b};
    \node[cell] at (60.4,\y) {\c};
  }
  \node[bar] at (36,3.5) {\iflangja{割り込み，ディスパッチャ}{interrupts, dispatcher}};
  \end{scope}
\end{tikzpicture}
   (width ~118 mm)
\begin{tikzpicture}[x=1mm, y=1mm, font=\scriptsize\sffamily,
  lay/.style={draw, rounded corners=1pt, fill=white, minimum width=40mm,
              text width=38mm, align=center, inner sep=0.5pt},
  bar/.style={draw, fill=ln-box, minimum width=72mm, minimum height=5mm,
              inner sep=0.5pt},
  hd/.style={draw, fill=ln-box, minimum width=23.2mm, minimum height=5mm,
             inner sep=0.5pt, font=\scriptsize\sffamily\bfseries},
  cell/.style={draw, fill=white, minimum width=23.2mm, text width=22.2mm,
               minimum height=9mm, align=center, inner sep=0.5pt},
  lab/.style={font=\scriptsize\sffamily, text=ln-muted},
  ttl/.style={font=\footnotesize\sffamily\bfseries, anchor=base}]
  % ---------------- layers
  \node[ttl] at (20,54) {\iflangja{Linux システムの階層}{layers of a Linux system}};
  \node[lay, minimum height=6mm] at (20,48.5) {\iflangja{ユーザ}{users}};
  \node[lay, minimum height=10mm] at (20,39)
    {\iflangja{標準ユーティリティ\\（シェル，エディタ，コンパイラ，…）}{standard utility programs\\(shell, editors, compilers, \dots)}};
  \node[lay, minimum height=10mm] at (20,27.5)
    {\iflangja{標準ライブラリ\\（open, close, read, write, fork, …）}{standard library\\(open, close, read, write, fork, \dots)}};
  \node[lay, minimum height=8mm, fill=ln-box, draw=ln-accent, thick] at (20,15.5)
    {\iflangja{Linux カーネル}{Linux kernel}};
  \node[lay, minimum height=8mm] at (20,5.5)
    {\iflangja{ハードウェア}{hardware}};
  \draw[dashed, ln-rule, thick] (-1,21) -- (41,21);
  % ---------------- kernel structure
  \begin{scope}[xshift=46mm]
  \node[ttl] at (36,54) {\iflangja{Linux カーネルの構成}{structure of the Linux kernel}};
  \node[bar] at (36,46) {\iflangja{システムコール}{system calls}};
  \node[hd] at (11.6,39.5) {I/O};
  \node[hd] at (36,39.5)   {\iflangja{メモリ管理}{memory mgmt.}};
  \node[hd] at (60.4,39.5) {\iflangja{プロセス管理}{process mgmt.}};
  \foreach \y/\a/\b/\c in {
      31.5/{\iflangja{端末\\文字デバイスドライバ}{terminals, character\\device drivers}}/{\iflangja{仮想メモリ}{virtual memory}}/{\iflangja{シグナル処理}{signal handling}},
      21.5/{\iflangja{ソケット，プロトコル，\\ネットワークドライバ}{sockets, protocols,\\network drivers}}/{\iflangja{ページング，\\ページ置換}{paging, page\\replacement}}/{\iflangja{プロセス・スレッドの\\生成と終了}{process and thread\\creation, termination}},
      11.5/{\iflangja{ファイルシステム，\\ブロックデバイスドライバ}{file systems, block\\device drivers}}/{\iflangja{ページキャッシュ}{page cache}}/{\iflangja{CPU スケジューリング}{CPU scheduling}}} {
    \node[cell] at (11.6,\y) {\a};
    \node[cell] at (36,\y)   {\b};
    \node[cell] at (60.4,\y) {\c};
  }
  \node[bar] at (36,3.5) {\iflangja{割り込み，ディスパッチャ}{interrupts, dispatcher}};
  \end{scope}
\end{tikzpicture}

  \caption{Left: the layers of a Linux system; the dashed line is the
    user/kernel boundary.  Right: the components of the Linux kernel
    (after Tanenbaum and Bos \cite{tanenbaum2015}).}
  \label{fig:l01-linux}
\end{figure}

The kernel itself (\cref{fig:l01-linux}, right) is entered through the
system calls at the top, and at the bottom it handles interrupts and
contains the dispatcher, which switches the CPU between processes.  In
between lie three components.  The I/O component contains the terminal
and character device drivers, the sockets, network protocols and network
device drivers, and the file systems and block device drivers.  The memory
management component implements virtual memory, paging and page
replacement, and the page cache.  The process management component handles
signals, creates and terminates processes and threads, and schedules the
CPU.  Each component can be modified or replaced independently of the
others, which is what makes the modular approach possible.

Mac OS X (now macOS) is built around the Mach microkernel
(\cref{fig:l01-macos}).\margin{Mach, BSD, the I/O Kit and kernel
extensions share one kernel address space: a hybrid design
(Silberschatz et al.).}  Mach implements the key
primitives: memory management, thread scheduling, and IPC, including
remote procedure calls (Lesson~13).  A BSD component provides Unix
interoperability through the BSD command-line interface, support for the
POSIX API and network I/O.  All application environments are layered above
these two.  Below them, the I/O Kit is an environment for developing device
drivers, and kernel extensions are components that are loaded into the
kernel dynamically.

\begin{figure}[htbp]
  \centering
  % Structure of Mac OS X: Mach microkernel with a BSD layer, I/O Kit and kernel extensions.
% I/O Kit and kernel extensions run in kernel mode, so they lie inside the
% kernel environment.
% Usage: % Structure of Mac OS X: Mach microkernel with a BSD layer, I/O Kit and kernel extensions.
% I/O Kit and kernel extensions run in kernel mode, so they lie inside the
% kernel environment.
% Usage: % Structure of Mac OS X: Mach microkernel with a BSD layer, I/O Kit and kernel extensions.
% I/O Kit and kernel extensions run in kernel mode, so they lie inside the
% kernel environment.
% Usage: \input{figures/l01-macos}   (width ~96 mm)
\begin{tikzpicture}[x=1mm, y=1mm, font=\scriptsize\sffamily,
  blk/.style={draw, rounded corners=1pt, fill=white, align=center, inner sep=1pt},
  lab/.style={font=\scriptsize\sffamily\bfseries, text=ln-accent}]
  \node[blk, minimum width=96mm, minimum height=7mm] at (48,43)
    {\iflangja{アプリケーション環境とサービス}{application environments and services}};
  \draw[draw=ln-accent, thick, fill=ln-box, rounded corners=2pt] (0,1) rectangle (96,38);
  \node[lab, anchor=north west] at (1,37.5) {\iflangja{カーネル環境}{kernel environment}};
  \node[blk, minimum width=92mm, minimum height=8mm] at (48,26)
    {\textbf{BSD}\quad \iflangja{BSD コマンドライン，POSIX API，ネットワーク I/O}{BSD command line, POSIX API, network I/O}};
  \node[blk, minimum width=92mm, minimum height=8mm] at (48,16)
    {\textbf{Mach}\quad \iflangja{メモリ管理，スレッドスケジューリング，IPC・RPC}{memory management, thread scheduling, IPC and RPC}};
  \node[blk, minimum width=45.5mm, minimum height=7mm] at (24.75,6.5)
    {\textbf{I/O Kit}\quad \iflangja{デバイスドライバ}{device drivers}};
  \node[blk, minimum width=45.5mm, minimum height=7mm] at (71.25,6.5)
    {\iflangja{カーネル拡張}{kernel extensions}};
\end{tikzpicture}
   (width ~96 mm)
\begin{tikzpicture}[x=1mm, y=1mm, font=\scriptsize\sffamily,
  blk/.style={draw, rounded corners=1pt, fill=white, align=center, inner sep=1pt},
  lab/.style={font=\scriptsize\sffamily\bfseries, text=ln-accent}]
  \node[blk, minimum width=96mm, minimum height=7mm] at (48,43)
    {\iflangja{アプリケーション環境とサービス}{application environments and services}};
  \draw[draw=ln-accent, thick, fill=ln-box, rounded corners=2pt] (0,1) rectangle (96,38);
  \node[lab, anchor=north west] at (1,37.5) {\iflangja{カーネル環境}{kernel environment}};
  \node[blk, minimum width=92mm, minimum height=8mm] at (48,26)
    {\textbf{BSD}\quad \iflangja{BSD コマンドライン，POSIX API，ネットワーク I/O}{BSD command line, POSIX API, network I/O}};
  \node[blk, minimum width=92mm, minimum height=8mm] at (48,16)
    {\textbf{Mach}\quad \iflangja{メモリ管理，スレッドスケジューリング，IPC・RPC}{memory management, thread scheduling, IPC and RPC}};
  \node[blk, minimum width=45.5mm, minimum height=7mm] at (24.75,6.5)
    {\textbf{I/O Kit}\quad \iflangja{デバイスドライバ}{device drivers}};
  \node[blk, minimum width=45.5mm, minimum height=7mm] at (71.25,6.5)
    {\iflangja{カーネル拡張}{kernel extensions}};
\end{tikzpicture}
   (width ~96 mm)
\begin{tikzpicture}[x=1mm, y=1mm, font=\scriptsize\sffamily,
  blk/.style={draw, rounded corners=1pt, fill=white, align=center, inner sep=1pt},
  lab/.style={font=\scriptsize\sffamily\bfseries, text=ln-accent}]
  \node[blk, minimum width=96mm, minimum height=7mm] at (48,43)
    {\iflangja{アプリケーション環境とサービス}{application environments and services}};
  \draw[draw=ln-accent, thick, fill=ln-box, rounded corners=2pt] (0,1) rectangle (96,38);
  \node[lab, anchor=north west] at (1,37.5) {\iflangja{カーネル環境}{kernel environment}};
  \node[blk, minimum width=92mm, minimum height=8mm] at (48,26)
    {\textbf{BSD}\quad \iflangja{BSD コマンドライン，POSIX API，ネットワーク I/O}{BSD command line, POSIX API, network I/O}};
  \node[blk, minimum width=92mm, minimum height=8mm] at (48,16)
    {\textbf{Mach}\quad \iflangja{メモリ管理，スレッドスケジューリング，IPC・RPC}{memory management, thread scheduling, IPC and RPC}};
  \node[blk, minimum width=45.5mm, minimum height=7mm] at (24.75,6.5)
    {\textbf{I/O Kit}\quad \iflangja{デバイスドライバ}{device drivers}};
  \node[blk, minimum width=45.5mm, minimum height=7mm] at (71.25,6.5)
    {\iflangja{カーネル拡張}{kernel extensions}};
\end{tikzpicture}

  \caption{The structure of Mac OS X: the Mach microkernel with a BSD
    layer, the I/O Kit and kernel extensions, all of which run in kernel
    mode.}
  \label{fig:l01-macos}
\end{figure}

\begin{keypoint}[Lesson summary]
  An operating system is the privileged layer of software between
  applications and hardware.  It hides hardware complexity behind
  abstractions, arbitrates the use of hardware resources, and isolates and
  protects applications from one another.  It is built from abstractions,
  mechanisms and policies; separating mechanism from policy and optimizing
  for the common case guide its design.  Hardware-supported user and kernel
  modes protect the operating system; applications cross the boundary
  through system calls and traps, which are not cheap, and the operating
  system notifies applications with signals.  Monolithic, modular and
  microkernel organizations trade completeness, flexibility, size and
  performance against each other, as the modular Linux kernel and the
  Mach-based Mac OS X illustrate.
\end{keypoint}

\end{document}
